\documentclass[a4paper,10pt]{article}
\usepackage[utf8]{inputenc}
\usepackage{graphicx}
\usepackage{subfigure}
\usepackage{fancyhdr}
\usepackage{listings}
\usepackage{ gensymb }


%opening

\pagestyle{fancy} %eigener Seitenstil
\fancyhf{} %alle Kopf- und Fußzeilenfelder bereinigen
\fancyhead[L]{SEP Abgabegespräch SS2015/16} %Kopfzeile links
\fancyhead[C]{} %zentrierte Kopfzeile
\fancyhead[R]{\includegraphics[width=1cm]{../../template/tugraz_bw}} %Kopfzeile rechts
\renewcommand{\headrulewidth}{0.4pt} %obere Trennlinie
\fancyfoot[C]{} %Seitennummer
\renewcommand{\footrulewidth}{0.4pt} %untere Trennlinie


\lstset{language=C++, basicstyle=\ttfamily\tiny}  

\begin{document}

\section*{Bank-Konto}

Entwickle eine einfache Konto-Software. Das Programm kann Konten erstellen, darauf Geld einzahlen und den Kontostand anzeigen. Dafür soll es 3 Befehle geben:

\begin{description}
 \item[create $<$kontoname$>$] Erstellt ein Konto mit dem angegebenen Namen und einem Betrag von 0.
 \item[add $<$kontoname$>$ $<$amount$>$] Fügt den Betrag dem angegebenen Konto hinzu. Es müssen keine Überprüfung gemacht werden, ob das Konto dabei negativ wird.
 \item[show $<$kontoname$>$] Zeigt den Namen und Kontostand des angegebenen Kontos an.
\end{description}

Sollte bei \texttt{add} oder \texttt{show} das Konto nicht existieren, soll ``Name does not exist!'' ausgegeben werden. Bei der Eingabe von \texttt{quit} oder EOF soll das Programm beendet werden.

\subsection*{Beispiel}
\begin{verbatim}
> create Peter
> show Peter
Peter: 0
> add Peter 10
> create Sam
> add Sam -5
> show Sam
Sam: -5
> show Peter
Peter: 10
> show Mike
Name does not exist!
\end{verbatim}

\subsection*{Aufruf}
\texttt{./konto}

\subsection*{Bewertung}
\begin{itemize}
 \item Dokumentation und Programmierstil: 1 Punkt
 \item Einhalten der OOP Konzepte: 1 Punkt
 \item Struktur \& Korrektheit: 1 Punkt
\end{itemize}

\newpage


\end{document}
